% =====================================================================
%  E8: every combination of five per-head filters (decoding, BOS, task
%  words, input, rest of the instruction). Record of the experiment: why,
%  design, the D-filter matching gap found after the first run and its fix,
%  sets, both checks, overlap with Davidson's FV heads.
%  Tables are generated by scripts/e8_tables.py (paper/tables/e8_*.tex).
% =====================================================================
\documentclass[10pt,a4paper]{article}

\usepackage[utf8]{inputenc}
\usepackage[T1]{fontenc}
\usepackage[margin=1.8cm]{geometry}
\usepackage{amsmath,amssymb}
\usepackage{booktabs,tabularx,array}
\usepackage{enumitem}
\usepackage{xcolor}
\usepackage{caption}
\usepackage{float}
\usepackage{graphicx}
\usepackage[colorlinks=true,linkcolor=blue!50!black,urlcolor=blue!60!black]{hyperref}

\setlist{nosep,leftmargin=*}
\setlength{\parindent}{0pt}
\setlength{\parskip}{5pt}
\renewcommand{\arraystretch}{1.12}
\captionsetup{font=small,labelfont=bf}

\title{E8: Every Combination of Five Head Filters\\\large Lexical task heads, simple tasks}
\author{Filippo Merlo}
\date{9 November 2026}

\begin{document}
\maketitle
\vspace{-1.5em}

\paragraph{In short.} E8 describes each head by five yes/no properties: it decodes the task's most active words under J-head (D), it does not look mainly at BOS (B), it attends to the instruction's task words (T), to the query (I), or not mainly to the rest of the instruction (R). Every combination was patched under E4's held-out protocol. Two sets stand out. Attention to the task words alone (T) is the best E8 set and the one chosen on the first fold. Decoding plus the BOS filter (DB) comes next with less than half as many heads; it needs no hand-written word list and no causal step. It is above CC-pooled and FV-20 and below FV-$N_t$. The two win different tasks: T the facts, DB the word relations, present--past and French. Attention to the query or to the rest of the instruction lowers every set it is added to. A first run had missed the country decoders: our lowercased word lists never matched the `` Capital'' tokens the lens decodes. With that fixed and the list lengthened to 30 words, every set that uses D improved, DB the most.

% ---------------------------------------------------------------------
\section{Why}

E7 showed that among heads that decode a task's words, the ones that look at the prompt rather than BOS carry the task. That leaves open which attention property matters: is ``not BOS'' the right filter, or do heads that attend to the task words, or to the item the task is applied to, do better? And does decoding add anything once attention is known? E8 tests every combination of one decoding filter and four attention filters, on the same items as E4 and E7.

% ---------------------------------------------------------------------
\section{Design}

E4's protocol, J-head only, rule style, the 14 simple tasks. All filters and thresholds were fixed before patching; the list behind D was changed after the first run (Sec.~\ref{sec:dfix}).

\begin{enumerate}
\item \textbf{Filters}, per task and head, on the choosing half's rule prompts (five wordings pooled, 500 prompts), at the last prompt token:
  \begin{itemize}
  \item \textbf{D} decoding: share of prompts whose J-head top 10 holds a token of the task's word list; pass if $\geq$0.10 (the CC / E7 head rule). List: the task's 30 most active decoded words, E7's word rules (top-10 tokens decoded, lowercased, alphabetic, $\geq$3 letters, no stop words, Zipf $\geq$3; activity = share of (prompt, head) pairs whose top 10 holds the word), nothing removed.
  \item \textbf{B} BOS: mean attention on BOS $\leq$0.5 (E7's filter).
  \item \textbf{T} task words: mean attention on the instruction's task-word tokens $\geq$0.05. Task words = Yang et al.'s term list plus the task's top-10 J-head ``specific'' words of E4 (choosing half), so T uses a hand-written list and the lens.
  \item \textbf{I} input: mean attention on the query $\geq$0.05.
  \item \textbf{R} rest: mean attention on the rest of the instruction $\leq$0.5.
  \end{itemize}
\item \textbf{Sets}: a set is named by its filters (DB = D and B); a head is in it if it passes all of them. All 31 non-empty combinations, all layers, no size cap; an empty set is left out for that task. No random controls.
\item \textbf{Patching} (E4): both folds; rule means of the choosing half and the control mean in both checks; $\alpha\in\{1,2\}$ only (to save compute; E7's chosen $\alpha$ were mostly 1--2); replies $\leq$16 tokens; flexible score; every answer saved per item. Check~1: items answered correctly with one shot, patched in their zero-shot prompts. Check~2: items answered wrongly with one shot, patched in their 1-shot prompts. Accuracy at the cross-fitted $\alpha$ (the best one on the other fold), mean of the two folds.
\item \textbf{References} on the same items and folds: E7's raw $+$F, E4's CC-20, CC, FV-20 and FV-$N_t$. Their runs had ten $\alpha$; here they are cross-fitted over $\alpha$ 1 and 2 only, so they are a little lower than in their own reports (CC 0.40 vs 0.44, FV-20 0.33 vs 0.38, raw $+$F 0.31 vs 0.40 in check~1).
\item \textbf{Filter choice without the reported items}: the set with the highest check-1 accuracy on fold~1 (each task at its own best $\alpha$) is picked; its fold-2 accuracy is the reported number.
\end{enumerate}

% ---------------------------------------------------------------------
\section{The decoding filter: a matching gap, and the list length}\label{sec:dfix}

\paragraph{First run.} D first used E7's raw list (the 10 most active words) and the word-to-token rule of Yang et al.'s code: each word is tried as written, with a leading space, in upper, capitalised and lower case, and after ``.'', ``\_'' or ``-''. D sets collapsed on landmark--country (0.02--0.03 in check~1). Looking at the heads behind it: L21H2 has `` England'' in its J-head top 10 on 488 of 500 prompts and \emph{england} is on the list, yet its D rate was 0. The rule never forms a leading space with a capital letter, which is how the lens decodes names. Yang et al.'s own lists are capitalised, so their code does not meet the problem; our data-driven lists are lowercased.

\paragraph{Fix.} For our lowercased lists, the leading-space capitalised and upper-case forms are added (`` Canada'', `` CANADA''; \texttt{lth.terms.expand\_lowercase\_to\_token\_ids}). The Yang reproduction keeps the exact rule.

\paragraph{List length.} Ten words is short: on landmark--country the list holds six country names, so heads decoding `` Japan'' or `` Italy'' are missed for a second reason. Before any rerun, list lengths of 10, 30 and 50 words were compared on the stored readings of the choosing halves, without patching, by how many of E4's FV-20 heads (chosen on the same half) end up in D and DB (Table~\ref{tab:listlen}).

\begin{table}[H]\centering\small
% Generated by scripts/e8_tables.py; do not edit.
\begin{tabular}{rlrrrrrrr}
\toprule
words & case fix & D & FV-20 in D & chance & DB & FV-20 in DB & chance & generic \\
\midrule
10 & no & 112 & 5.2 & 2.2 & 32 & 4.7 & 0.6 & 4/10 \\
10 & yes & 137 & 6.9 & 2.7 & 41 & 6.0 & 0.8 & 4/10 \\
30 & no & 213 & 8.1 & 4.2 & 56 & 7.1 & 1.1 & 13/30 \\
30 & yes & 245 & 10.0 & 4.8 & 65 & 8.5 & 1.3 & 13/30 \\
50 & no & 273 & 9.3 & 5.3 & 70 & 7.8 & 1.4 & 23/50 \\
50 & yes & 310 & 11.2 & 6.1 & 80 & 9.3 & 1.6 & 23/50 \\
\bottomrule
\end{tabular}

\caption{D-filter variants, mean over tasks and folds. Chance: FV-20 heads expected in a random set of the same size. Generic: list words among the 50 words heads decode most on WikiText. Source: \texttt{scripts/e8\_list\_length.py}.}
\label{tab:listlen}
\end{table}

The case fix adds FV heads at every length without lowering the ratio to chance. A longer list adds more, but less selectively: at 50 words almost half the list is generic and D holds 30\% of all heads; from 30 to 50 words DB grows by 15 heads for about one more FV head. We chose 30 words with the fix. This choice was prompted by a patched result (the collapse), but made without patched accuracies. The 15 sets without D were unchanged; the 16 with D were rebuilt and rerun on the same items.

\begin{table}[H]\centering\small
% Generated by scripts/e8_tables.py; do not edit.
\begin{tabular}{lrrrrrr}
\toprule
set & heads, 1st & heads, now & check 1, 1st & check 1, now & check 2, 1st & check 2, now \\
\midrule
DB & 32 & 65 & 0.31 & 0.51 & 0.34 & 0.32 \\
DT & 34 & 64 & 0.27 & 0.49 & 0.32 & 0.35 \\
DBT & 20 & 38 & 0.23 & 0.49 & 0.32 & 0.35 \\
D & 110 & 249 & 0.34 & 0.41 & 0.32 & 0.26 \\
DTR & 30 & 59 & 0.18 & 0.27 & 0.28 & 0.30 \\
DBTR & 16 & 30 & 0.17 & 0.25 & 0.27 & 0.31 \\
DBI & 21 & 43 & 0.09 & 0.25 & 0.30 & 0.24 \\
DBR & 23 & 46 & 0.13 & 0.22 & 0.28 & 0.28 \\
DI & 48 & 107 & 0.12 & 0.22 & 0.29 & 0.26 \\
DTI & 20 & 38 & 0.08 & 0.21 & 0.28 & 0.29 \\
DR & 103 & 229 & 0.16 & 0.20 & 0.28 & 0.25 \\
DBTI & 13 & 23 & 0.06 & 0.18 & 0.27 & 0.30 \\
DTIR & 18 & 34 & 0.06 & 0.13 & 0.25 & 0.27 \\
DIR & 44 & 98 & 0.06 & 0.09 & 0.26 & 0.23 \\
DBTIR & 11 & 20 & 0.05 & 0.09 & 0.23 & 0.27 \\
DBIR & 16 & 35 & 0.03 & 0.08 & 0.23 & 0.22 \\
\bottomrule
\end{tabular}

\caption{Every set with D, first run (10 words, Yang's expansion) and now (30 words, with the fix): heads (median over tasks), task accuracy at the cross-fitted $\alpha$, mean over tasks.}
\label{tab:firstrun}
\end{table}

All 16 sets rise in check~1, those with B or T the most; sets roughly double in size. In check~2 the changes are small and go both ways.

% ---------------------------------------------------------------------
\section{Word lists and sets}

\begin{table}[H]\centering\scriptsize
% Generated by scripts/e8_tables.py; do not edit.
\begin{tabular}{>{\raggedright\arraybackslash}p{2.2cm}>{\raggedright\arraybackslash}p{12.4cm}r}
\toprule
Task & 30 most active decoded words (fold 1, most active first) & fold 2 same \\
\midrule
antonym & \emph{opposite, what, which, word, instead, etc, because, http, opposing, words, therefore, amongst, definition, colour, reverse, synonyms, https, why, mph, maybe, only, whats, name, their, itself, they, meaning, american, ant, answer} & 28 \\
synonym & \emph{what, which, etc, synonyms, amongst, https, word, maybe, whilst, their, words, colour, similar, synonym, http, definition, name, whats, english, then, syn, they, why, same, mph, afterward, possibly, herself, her, being} & 28 \\
singular--plural & \emph{because, what, single, which, https, http, name, noun, etc, multiple, plural, mph, maybe, amongst, either, linguistic, therefore, names, named, then, naming, whilst, even, only, colour, similar, english, called, there, whats} & 30 \\
present--past & \emph{what, etc, english, https, verb, http, which, grammar, name, maybe, because, linguistic, really, question, verbs, they, why, colour, mph, even, then, amongst, their, called, word, how, only, britain, whats, british} & 29 \\
country--capital & \emph{city, cities, which, country, what, town, countries, answer, mph, international, only, maybe, http, https, geography, downtown, colour, capital, canada, really, american, there, itself, they, pretty, answers, italy, london, etc, our} & 28 \\
country--currency & \emph{currency, currencies, country, which, monetary, money, what, countries, dollar, international, http, https, financial, maybe, american, answer, amongst, foreign, colour, whilst, mph, they, germany, really, etc, british, even, canada, only, also} & 28 \\
landmark--country & \emph{country, which, countries, england, canada, germany, what, france, location, australia, city, international, geography, local, italy, located, mph, http, maybe, geographical, locations, etc, https, usa, county, geographic, being, cities, netherlands, amongst} & 29 \\
park--country & \emph{country, which, countries, what, england, canada, location, australia, park, located, france, national, mph, tourism, germany, geography, unesco, parks, international, italy, city, nation, etc, tourist, geographical, county, itself, also, wildlife, locations} & 30 \\
product--company & \emph{company, companies, which, manufacturer, what, product, http, etc, maybe, microsoft, mph, https, manufacturers, products, because, only, ones, really, who, itself, named, name, even, why, colour, owned, wikipedia, software, apple, ownership} & 29 \\
next-item & \emph{next, which, what, january, july, september, november, october, february, mph, maybe, sequence, because, etc, question, items, december, june, after, item, april, herself, pretty, even, calendar, then, answer, amongst, christmas, really} & 24 \\
prev-item & \emph{previous, what, preceding, prior, answer, question, item, which, next, because, mph, why, etc, items, amongst, sequence, january, previously, http, even, earlier, september, maybe, answers, https, how, november, october, she, whilst} & 25 \\
English--French & \emph{french, france, english, paris, what, which, translation, http, etc, their, italian, european, american, they, translations, language, maybe, https, linguistic, then, japanese, german, being, colour, called, bilingual, therefore, languages, how, named} & 29 \\
English--German & \emph{german, germany, germans, english, which, http, what, berlin, european, french, etc, american, translation, linguistic, language, https, maybe, their, called, colour, therefore, italian, von, japanese, munich, czech, translations, languages, being, die} & 30 \\
English--Spanish & \emph{spanish, english, what, mexican, http, hispanic, translation, bilingual, spain, which, language, etc, mexico, their, translations, https, portuguese, linguistic, french, called, languages, american, being, name, they, maybe, madrid, santiago, why, therefore} & 28 \\
\bottomrule
\end{tabular}

\caption{D's word lists, fold 1; last column: words also on fold 2's list.}
\label{tab:lists}
\end{table}

The folds agree on 24--30 of the 30 words. Between 6 and 18 words per list are generic (\emph{what, which, http, mph, maybe}); the rest name the task or its answers (country names on landmark and park, months on next-item, \emph{opposite} on antonym).

\begin{table}[H]\centering\small
\resizebox{\linewidth}{!}{% Generated by scripts/e8_tables.py; do not edit.
\begin{tabular}{lrrrrrrrrrrrr}
\toprule
Task & D & B & T & I & R & DB & DT & BT & DBT & DBTIR & D (1st) & DB (1st) \\
\midrule
antonym & 241 & 192 & 206 & 360 & 970 & 67 & 69 & 110 & 44 & 26 & 124 & 42 \\
synonym & 246 & 167 & 156 & 322 & 971 & 56 & 58 & 93 & 38 & 22 & 134 & 34 \\
singular--plural & 230 & 210 & 124 & 332 & 948 & 67 & 48 & 79 & 32 & 19 & 119 & 38 \\
present--past & 221 & 200 & 13 & 337 & 916 & 52 & 3 & 12 & 3 & 0 & 108 & 28 \\
country--capital & 269 & 174 & 200 & 320 & 968 & 64 & 84 & 95 & 44 & 28 & 117 & 32 \\
country--currency & 262 & 162 & 149 & 254 & 960 & 52 & 61 & 74 & 32 & 18 & 106 & 27 \\
landmark--country & 264 & 212 & 85 & 364 & 944 & 84 & 43 & 47 & 24 & 15 & 105 & 32 \\
park--country & 210 & 210 & 113 & 451 & 958 & 68 & 45 & 58 & 24 & 16 & 116 & 36 \\
product--company & 274 & 199 & 57 & 442 & 958 & 82 & 30 & 40 & 21 & 14 & 132 & 35 \\
next-item & 264 & 160 & 302 & 359 & 990 & 64 & 99 & 104 & 47 & 38 & 106 & 25 \\
prev-item & 252 & 224 & 150 & 288 & 924 & 74 & 66 & 98 & 38 & 28 & 106 & 36 \\
English--French & 254 & 193 & 212 & 326 & 952 & 66 & 74 & 107 & 41 & 22 & 112 & 28 \\
English--German & 224 & 194 & 218 & 310 & 948 & 58 & 74 & 112 & 40 & 21 & 80 & 18 \\
English--Spanish & 222 & 190 & 230 & 286 & 957 & 62 & 77 & 114 & 44 & 19 & 95 & 30 \\
\midrule
\textbf{median} & 249 & 194 & 154 & 329 & 958 & 65 & 64 & 94 & 38 & 20 & 110 & 32 \\
\bottomrule
\end{tabular}
}
\caption{Heads per set, mean of the two folds (1{,}024 heads in all). D (1st), DB (1st): the first run's D and DB.}
\label{tab:sizes}
\end{table}

R passes nearly every head and I about a third, so neither selects much on its own. T varies most across tasks: 13 heads on present--past, about 300 on next-item. DB keeps a quarter to a third of D on every task (52--84 heads).

% ---------------------------------------------------------------------
\section{Results}

\subsection{Check 1}

\begin{table}[H]\centering\small
\resizebox{\linewidth}{!}{% Generated by scripts/e8_tables.py; do not edit.
\begin{tabular}{lrrrrrrrrrrrrrr}
\toprule
Task & items & unp. & D & B & T & DB & DT & BT & DBT & E7 raw+F & CC-20 & CC & FV-20 & FV-$N_t$ \\
\midrule
antonym & 100/100 & 0.00 & \textbf{0.74} & 0.00 & 0.60 & 0.70 & 0.65 & 0.34 & 0.61 & 0.46 & 0.00 & 0.42 & 0.02 & 0.49 \\
synonym & 100/100 & 0.03 & 0.49 & 0.01 & 0.18 & 0.38 & 0.32 & 0.14 & 0.27 & 0.30 & 0.09 & 0.34 & 0.18 & \textbf{0.55} \\
singular--plural & 91/88 & 0.00 & 0.80 & 0.02 & 0.80 & 0.78 & 0.70 & 0.82 & 0.80 & \textbf{0.88} & 0.09 & 0.69 & 0.44 & 0.85 \\
present--past & 100/100 & 0.00 & 0.44 & 0.02 & 0.00 & 0.54 & 0.00 & 0.00 & 0.00 & 0.17 & 0.00 & 0.04 & 0.64 & \textbf{0.85} \\
country--capital & 84/86 & 0.00 & 0.65 & 0.46 & \textbf{0.98} & 0.69 & 0.89 & 0.93 & 0.94 & 0.87 & 0.77 & 0.75 & 0.88 & 0.94 \\
country--currency & 67/65 & 0.00 & 0.28 & 0.40 & \textbf{0.61} & 0.45 & 0.30 & 0.58 & 0.27 & 0.32 & 0.14 & 0.29 & 0.00 & 0.47 \\
landmark--country & 100/100 & 0.01 & 0.23 & 0.77 & \textbf{0.93} & 0.73 & 0.89 & 0.88 & 0.81 & 0.03 & 0.53 & 0.92 & 0.92 & 0.93 \\
park--country & 100/100 & 0.00 & 0.47 & 0.17 & 0.67 & 0.47 & 0.54 & 0.63 & 0.55 & 0.42 & 0.40 & 0.73 & 0.82 & \textbf{0.88} \\
product--company & 100/100 & 0.07 & 0.21 & 0.51 & 0.82 & 0.40 & 0.76 & 0.83 & 0.74 & 0.37 & 0.27 & 0.24 & 0.64 & \textbf{0.91} \\
next-item & 58/63 & 0.00 & 0.01 & 0.05 & 0.10 & 0.03 & 0.11 & 0.07 & 0.10 & 0.02 & 0.00 & 0.17 & 0.01 & \textbf{0.70} \\
prev-item & 53/50 & 0.00 & 0.01 & 0.02 & 0.04 & 0.04 & 0.02 & 0.08 & 0.03 & 0.16 & 0.00 & 0.13 & 0.01 & \textbf{0.19} \\
English--French & 100/100 & 0.00 & 0.34 & 0.07 & 0.49 & \textbf{0.57} & 0.47 & 0.20 & 0.49 & 0.13 & 0.00 & 0.26 & 0.00 & 0.21 \\
English--German & 100/100 & 0.00 & 0.46 & 0.14 & \textbf{0.67} & 0.60 & 0.57 & 0.43 & 0.50 & 0.07 & 0.00 & 0.18 & 0.00 & 0.02 \\
English--Spanish & 100/100 & 0.00 & 0.55 & 0.10 & 0.76 & 0.71 & 0.70 & 0.38 & \textbf{0.77} & 0.17 & 0.00 & 0.38 & 0.01 & 0.25 \\
\midrule
\textbf{mean} &  & 0.01 & 0.41 & 0.20 & 0.55 & 0.51 & 0.49 & 0.45 & 0.49 & 0.31 & 0.16 & 0.40 & 0.33 & \textbf{0.59} \\
\midrule
control mean &  &  & 0.02 & 0.00 & 0.02 & 0.01 & 0.03 & 0.02 & 0.03 & 0.02 & 0.03 & 0.05 & 0.03 & 0.04 \\
\bottomrule
\end{tabular}
}
\caption{Check 1, task accuracy at the cross-fitted $\alpha$, mean of the two folds; the main E8 sets and the references (E7 raw $+$F; E4's CC-20, CC, FV-20, FV-$N_t$), all cross-fitted over $\alpha$ 1 and 2. Items: patched items in fold 1/2. Control mean: accuracy with the control task's mean. All 31 sets per task: Table~\ref{tab:all1}.}
\label{tab:check1}
\end{table}

\begin{table}[H]\centering\small
% Generated by scripts/e8_tables.py; do not edit.
\begin{tabular}{lrrrrrr}
\toprule
set & check 1 & ctrl & check 2 & ctrl & heads (median) & tasks \\
\midrule
T & 0.55 & 0.02 & 0.40 & 0.11 & 154 & 14 \\
DB & 0.51 & 0.01 & 0.32 & 0.09 & 65 & 14 \\
DT & 0.49 & 0.03 & 0.35 & 0.12 & 64 & 14 \\
DBT & 0.49 & 0.03 & 0.35 & 0.14 & 38 & 14 \\
BT & 0.45 & 0.02 & 0.38 & 0.13 & 94 & 14 \\
D & 0.41 & 0.02 & 0.26 & 0.05 & 249 & 14 \\
TR & 0.31 & 0.01 & 0.35 & 0.10 & 137 & 14 \\
DTR & 0.27 & 0.04 & 0.30 & 0.11 & 59 & 14 \\
DBTR & 0.25 & 0.03 & 0.31 & 0.12 & 30 & 14 \\
DBI & 0.25 & 0.01 & 0.24 & 0.11 & 43 & 14 \\
DBR & 0.22 & 0.01 & 0.28 & 0.08 & 46 & 14 \\
BTR & 0.22 & 0.01 & 0.32 & 0.10 & 74 & 14 \\
DI & 0.22 & 0.01 & 0.26 & 0.12 & 107 & 14 \\
DTI & 0.21 & 0.03 & 0.29 & 0.15 & 38 & 14 \\
DR & 0.20 & 0.01 & 0.25 & 0.05 & 229 & 14 \\
B & 0.20 & 0.00 & 0.17 & 0.08 & 194 & 14 \\
DBTI & 0.18 & 0.03 & 0.30 & 0.15 & 23 & 14 \\
TI & 0.15 & 0.01 & 0.27 & 0.13 & 88 & 14 \\
BTI & 0.13 & 0.01 & 0.26 & 0.12 & 58 & 14 \\
DTIR & 0.13 & 0.04 & 0.27 & 0.15 & 34 & 13 \\
R & 0.12 & 0.00 & 0.06 & 0.01 & 958 & 14 \\
DIR & 0.09 & 0.01 & 0.23 & 0.11 & 98 & 14 \\
DBTIR & 0.09 & 0.03 & 0.27 & 0.14 & 20 & 13 \\
DBIR & 0.08 & 0.01 & 0.22 & 0.10 & 35 & 14 \\
TIR & 0.06 & 0.01 & 0.24 & 0.11 & 80 & 14 \\
BR & 0.06 & 0.00 & 0.17 & 0.07 & 124 & 14 \\
I & 0.05 & 0.01 & 0.06 & 0.02 & 329 & 14 \\
BTIR & 0.05 & 0.01 & 0.22 & 0.10 & 49 & 14 \\
BI & 0.04 & 0.00 & 0.13 & 0.09 & 120 & 14 \\
IR & 0.03 & 0.00 & 0.07 & 0.02 & 298 & 14 \\
BIR & 0.00 & 0.00 & 0.12 & 0.09 & 92 & 14 \\
\midrule
E7 raw+F & 0.31 & 0.02 & 0.34 & 0.13 &  &  \\
CC-20 & 0.16 & 0.03 & 0.27 & 0.10 &  &  \\
CC & 0.40 & 0.05 & 0.34 & 0.13 &  &  \\
FV-20 & 0.33 & 0.03 & 0.34 & 0.10 &  &  \\
FV-$N_t$ & 0.59 & 0.04 & 0.39 & 0.12 &  &  \\
\bottomrule
\end{tabular}

\caption{Every set, mean over the 14 tasks, task and control mean in both checks, sorted by check 1; heads: median over tasks; tasks: tasks with a non-empty set.}
\label{tab:overview}
\end{table}

\begin{itemize}
\item \textbf{T and DB lead, on different tasks.} T carries the fact tasks (capitals, currency, landmark, park, product) best among the E8 sets and leads DB on German and Spanish; DB carries antonym, synonym, present--past and English--French better. On present--past T has only 13 heads and does nothing; DB carries it, though less than FV-$N_t$.
\item \textbf{DB is the best set without a hand-written word list.} With about half CC's heads it is above CC on 9 of 14 tasks and above FV-20 on 9; it stays below FV-$N_t$ on 10 of 14.
\item \textbf{The BOS filter helps only decoders.} It raises D by 0.10 but lowers T by 0.09; on the other sets it changes little or lowers accuracy. Heads chosen by what they decode need it; heads chosen by where they look do not.
\item \textbf{T helps most combinations} (12 of 15). Added to DB (DBT, 38 heads), it keeps most of T's fact-task accuracy with a quarter of T's heads, but loses present--past.
\item \textbf{Input and rest attention hurt.} Adding I or R lowers the mean of every one of the 15 sets it is added to. Alone they do almost nothing.
\item \textbf{Sequences fail on every E8 set} (next- and prev-item $\leq$0.11); only FV-$N_t$ carries next-item.
\item Control means stay low on average. T, BT, DT, DBT and some of their extensions have a high control on park--country (0.18--0.34), so part of their park effect is not task-specific.
\end{itemize}

\begin{table}[H]\centering\small
% Generated by scripts/e8_tables.py; do not edit.
\begin{tabular}{lrr}
\toprule
set & fold 1, own best $\alpha$ & fold 2, cross-fitted \\
\midrule
T & 0.55 & 0.55 \\
DB & 0.53 & 0.49 \\
DT & 0.50 & 0.49 \\
DBT & 0.50 & 0.48 \\
BT & 0.46 & 0.44 \\
D & 0.41 & 0.41 \\
TR & 0.30 & 0.31 \\
DTR & 0.29 & 0.26 \\
\bottomrule
\end{tabular}

\caption{Filter choice: check-1 accuracy on fold 1 at each task's own best $\alpha$ (top 8 sets), and the same set on fold 2 at the cross-fitted $\alpha$, mean over tasks.}
\label{tab:select}
\end{table}

Chosen on fold 1, T keeps its accuracy on fold 2. DB, DT and DBT lose a little and end level with each other.

\subsection{Check 2}

\begin{table}[H]\centering\small
\resizebox{\linewidth}{!}{% Generated by scripts/e8_tables.py; do not edit.
\begin{tabular}{lrrrrrrrrrrrrrr}
\toprule
Task & items & unp. & D & B & T & DB & DT & BT & DBT & E7 raw+F & CC-20 & CC & FV-20 & FV-$N_t$ \\
\midrule
antonym & 100/100 & 0.00 & 0.29 & 0.04 & 0.37 & 0.33 & \textbf{0.39} & 0.35 & 0.36 & 0.34 & 0.31 & 0.35 & 0.30 & 0.35 \\
synonym & 100/100 & 0.00 & 0.11 & 0.02 & 0.07 & 0.16 & 0.16 & 0.10 & 0.14 & 0.16 & 0.13 & 0.17 & 0.11 & \textbf{0.18} \\
singular--plural & 11/15 & 0.00 & 0.73 & 0.21 & 0.76 & \textbf{0.79} & 0.76 & 0.76 & 0.75 & 0.77 & 0.68 & 0.77 & 0.64 & 0.65 \\
present--past & 30/31 & 0.00 & 0.68 & 0.12 & 0.75 & 0.69 & 0.46 & 0.74 & 0.46 & 0.78 & 0.57 & 0.64 & \textbf{0.89} & 0.82 \\
country--capital & 14/13 & 0.00 & 0.52 & 0.27 & \textbf{0.81} & 0.68 & \textbf{0.81} & 0.76 & 0.71 & 0.74 & 0.73 & 0.77 & 0.67 & 0.74 \\
country--currency & 31/34 & 0.00 & \textbf{0.37} & 0.19 & 0.29 & 0.22 & 0.17 & 0.25 & 0.16 & 0.12 & 0.12 & 0.23 & 0.19 & 0.17 \\
landmark--country & 100/87 & 0.00 & 0.04 & 0.19 & 0.31 & 0.22 & 0.33 & 0.33 & \textbf{0.34} & 0.19 & 0.32 & 0.32 & 0.33 & 0.32 \\
park--country & 22/23 & 0.00 & 0.14 & \textbf{0.45} & 0.38 & 0.22 & 0.26 & 0.32 & 0.26 & 0.38 & 0.21 & 0.33 & 0.33 & 0.38 \\
product--company & 84/75 & 0.00 & 0.11 & 0.18 & \textbf{0.30} & 0.16 & 0.23 & 0.28 & 0.22 & 0.21 & 0.13 & 0.16 & 0.20 & 0.22 \\
next-item & 54/50 & 0.00 & 0.02 & 0.11 & 0.44 & 0.12 & 0.35 & 0.50 & 0.55 & 0.25 & 0.08 & 0.34 & 0.23 & \textbf{0.61} \\
prev-item & 59/63 & 0.00 & 0.03 & 0.09 & 0.17 & 0.05 & 0.10 & 0.19 & 0.10 & 0.12 & 0.09 & 0.09 & 0.17 & \textbf{0.32} \\
English--French & 100/100 & 0.00 & 0.17 & 0.15 & 0.35 & 0.25 & 0.32 & 0.22 & 0.29 & 0.22 & 0.12 & 0.16 & 0.26 & \textbf{0.36} \\
English--German & 100/100 & 0.00 & 0.27 & 0.15 & \textbf{0.30} & 0.27 & 0.27 & 0.24 & 0.24 & 0.27 & 0.13 & 0.21 & 0.21 & 0.17 \\
English--Spanish & 100/100 & 0.00 & 0.21 & 0.24 & \textbf{0.35} & 0.29 & 0.33 & 0.28 & 0.30 & 0.25 & 0.14 & 0.24 & 0.20 & 0.22 \\
\midrule
\textbf{mean} &  & 0.00 & 0.26 & 0.17 & \textbf{0.40} & 0.32 & 0.35 & 0.38 & 0.35 & 0.34 & 0.27 & 0.34 & 0.34 & 0.39 \\
\midrule
control mean &  &  & 0.05 & 0.08 & 0.11 & 0.09 & 0.12 & 0.13 & 0.14 & 0.13 & 0.10 & 0.13 & 0.10 & 0.12 \\
\bottomrule
\end{tabular}
}
\caption{Check 2 (1-shot prompts answered wrongly), task accuracy at the cross-fitted $\alpha$, as Table~\ref{tab:check1}. All 31 sets per task: Table~\ref{tab:all2}.}
\label{tab:check2}
\end{table}

As in E4 and E7, the sets are close. T is the best E8 set and level with FV-$N_t$; DB is a little below CC. Sets with T gain most on next-item, but their control means there are close to the task means, so that gain is not task-specific. Several tasks have few items in this check (Table~\ref{tab:check2}), so single-task differences are coarse.

% ---------------------------------------------------------------------
\section{The E8 heads and the FV heads}

\begin{table}[H]\centering\small
\resizebox{\linewidth}{!}{% Generated by scripts/e8_tables.py; do not edit.
\begin{tabular}{lrrrrrrrr}
\toprule
Task & D/FV-20 & D/FV-$N_t$ & T/FV-20 & T/FV-$N_t$ & DB/FV-20 & DB/FV-$N_t$ & DBT/FV-20 & DBT/FV-$N_t$ \\
\midrule
antonym & 13.0 (4.7) & 50.0 (27.9) & 13.0 (4.0) & 59.5 (23.8) & 8.5 (1.3) & 33.5 (7.7) & 7.0 (0.8) & 22.0 (5.0) \\
synonym & 11.5 (4.8) & 43.5 (28.0) & 14.5 (3.1) & 49.5 (17.8) & 10.0 (1.1) & 26.0 (6.4) & 9.0 (0.8) & 18.0 (4.4) \\
singular--plural & 12.0 (4.5) & 48.0 (26.7) & 14.0 (2.4) & 45.0 (14.3) & 12.0 (1.3) & 34.0 (7.8) & 11.0 (0.6) & 20.5 (3.7) \\
present--past & 7.5 (4.3) & 31.5 (22.0) & 3.0 (0.3) & 8.0 (1.3) & 6.5 (1.0) & 19.0 (5.2) & 2.0 (0.1) & 3.0 (0.3) \\
country--capital & 14.5 (5.3) & 60.5 (34.1) & 16.0 (3.9) & 58.0 (25.4) & 11.5 (1.2) & 27.5 (8.1) & 10.0 (0.8) & 21.0 (5.5) \\
country--currency & 4.5 (5.1) & 36.0 (29.5) & 8.0 (2.9) & 36.5 (16.7) & 4.0 (1.0) & 17.0 (5.8) & 2.5 (0.6) & 11.5 (3.6) \\
landmark--country & 16.0 (5.2) & 68.0 (33.5) & 9.0 (1.7) & 29.5 (10.8) & 15.0 (1.6) & 43.5 (10.7) & 7.0 (0.5) & 16.0 (3.0) \\
park--country & 15.5 (4.1) & 63.5 (30.1) & 12.5 (2.2) & 37.0 (16.2) & 12.5 (1.3) & 36.5 (9.7) & 8.5 (0.5) & 14.5 (3.5) \\
product--company & 14.5 (5.3) & 55.5 (33.9) & 13.0 (1.1) & 28.0 (7.1) & 14.0 (1.6) & 37.0 (10.2) & 10.0 (0.4) & 16.5 (2.6) \\
next-item & 9.0 (5.1) & 51.0 (32.9) & 12.5 (5.9) & 71.0 (37.7) & 6.5 (1.3) & 24.5 (8.0) & 6.0 (0.9) & 22.0 (5.8) \\
prev-item & 5.5 (4.9) & 35.5 (27.1) & 9.0 (2.9) & 37.0 (16.2) & 5.0 (1.4) & 23.5 (7.9) & 3.0 (0.7) & 16.0 (4.1) \\
English--French & 7.0 (5.0) & 42.5 (29.6) & 6.0 (4.1) & 41.5 (24.7) & 6.0 (1.3) & 23.0 (7.6) & 4.5 (0.8) & 15.0 (4.8) \\
English--German & 4.5 (4.4) & 21.5 (21.5) & 9.5 (4.2) & 35.0 (20.8) & 4.0 (1.1) & 13.5 (5.5) & 3.0 (0.8) & 8.5 (3.8) \\
English--Spanish & 4.5 (4.3) & 30.0 (24.3) & 7.5 (4.5) & 50.5 (25.2) & 3.0 (1.2) & 15.5 (6.8) & 3.0 (0.9) & 12.0 (4.9) \\
\midrule
\textbf{mean} & 10.0 (4.8) & 45.5 (28.7) & 10.5 (3.1) & 41.9 (18.4) & 8.5 (1.3) & 26.7 (7.7) & 6.2 (0.7) & 15.5 (3.9) \\
\bottomrule
\end{tabular}
}
\caption{E4's FV-20 and FV-$N_t$ heads (FV-$N_t$: as many as CC, about 120) found in D, T, DB and DBT, per task, mean of the two folds; in brackets the count expected for a random set of the same size.}
\label{tab:overlap}
\end{table}

DB holds about 6 times as many FV-20 heads as a random set of its size, T about 3 times. DB finds most of FV-20 on the fact tasks except currency, and on singular--plural. On the translations and currency it holds few FV heads, though it carries the translations; there the two selections pick different heads.

% ---------------------------------------------------------------------
\section{Takeaways}

\begin{enumerate}
\item Two cheap selections each come close to the best causally chosen set: \textbf{attention to the task words (T)}, and \textbf{decoding the task's words while not looking at BOS (DB)}. They carry different tasks.
\item DB uses only the lens and attention, with no hand-written list or causal step, and is above CC and FV-20 on most tasks.
\item Attention to the query or to the rest of the instruction does not pick task heads.
\item How decoded tokens are matched to a word list matters: the leading-space capitalised forms hid most country decoders.
\end{enumerate}

\paragraph{Caveats.} One run per setting, greedy, no intervals, no random controls in E8. Set sizes differ (DB about 65 heads, T about 150, FV-20 20, FV-$N_t$ about 120), and earlier work found that size alone moves accuracy; the comparisons are not size-matched. D's list was changed after the first run (Sec.~\ref{sec:dfix}). Only $\alpha$ 1 and 2 were tried. T's task-word positions come from a list that includes Yang et al.'s terms. E7's and E4's CC lists still use Yang's expansion and have the same matching gap. Some tasks have few patched items in check~2.

\paragraph{Open next steps.}
\begin{itemize}
\item Size-matched controls for DB and T: random-drop and layer-matched draws (as E7), and FV sets of exactly DB's size.
\item Rebuild E7's raw $+$F and E4's CC with the case fix, to see how much of their landmark collapse was the gap.
\item The union of DB and T, since they win different tasks.
\item Why no attention or decoding filter finds the sequence heads that FV-$N_t$ uses on next-item.
\end{itemize}

\paragraph{Reproduce.}
\begin{itemize}
\item \texttt{scripts/e8\_list\_length.py}: the list-length check (Table~\ref{tab:listlen}).
\item \texttt{scripts/e8\_sets.py --choose A|B}: word lists and sets of each fold.
\item \texttt{scripts/e8\_patch.sh}: patching (GPU, about 45 min on 4 GPUs for all 31 sets); the D rerun with \texttt{SETS="D DB \dots" PART=d30}.
\item \texttt{scripts/e8\_report.py}: accuracies, \texttt{results/e8\_filters/report.md}.
\item \texttt{scripts/e8\_tables.py}: these tables.
\item First run: commit \texttt{4d19133}; its summaries in \texttt{results/\allowbreak e8\_filters/\allowbreak first\_run\_d10}.
\end{itemize}

\appendix
\section{All sets, per task}

\begin{table}[H]\centering\scriptsize
\resizebox{\linewidth}{!}{% Generated by scripts/e8_tables.py; do not edit.
\begin{tabular}{lrrrrrrrrrrrrrrr}
\toprule
set & ant. & syn. & sg--pl & pr--pa & capital & currency & landm. & park & product & next & prev & En--Fr & En--De & En--Es & mean \\
\midrule
D & 0.74 & 0.49 & 0.80 & 0.44 & 0.65 & 0.28 & 0.23 & 0.47 & 0.21 & 0.01 & 0.01 & 0.34 & 0.46 & 0.55 & 0.41 \\
B & 0.00 & 0.01 & 0.02 & 0.02 & 0.46 & 0.40 & 0.77 & 0.17 & 0.51 & 0.05 & 0.02 & 0.07 & 0.14 & 0.10 & 0.20 \\
T & 0.60 & 0.18 & 0.80 & 0.00 & 0.98 & 0.61 & 0.93 & 0.67 & 0.82 & 0.10 & 0.04 & 0.49 & 0.67 & 0.76 & 0.55 \\
I & 0.00 & 0.01 & 0.00 & 0.00 & 0.00 & 0.00 & 0.17 & 0.13 & 0.28 & 0.02 & 0.00 & 0.09 & 0.04 & 0.01 & 0.05 \\
R & 0.00 & 0.01 & 0.01 & 0.01 & 0.10 & 0.40 & 0.42 & 0.01 & 0.21 & 0.02 & 0.01 & 0.15 & 0.07 & 0.24 & 0.12 \\
\midrule
DB & 0.70 & 0.38 & 0.78 & 0.54 & 0.69 & 0.45 & 0.73 & 0.47 & 0.40 & 0.03 & 0.04 & 0.57 & 0.60 & 0.71 & 0.51 \\
DT & 0.65 & 0.32 & 0.70 & 0.00 & 0.89 & 0.30 & 0.89 & 0.54 & 0.76 & 0.11 & 0.02 & 0.47 & 0.57 & 0.70 & 0.49 \\
DI & 0.37 & 0.17 & 0.25 & 0.36 & 0.37 & 0.03 & 0.17 & 0.33 & 0.23 & 0.00 & 0.00 & 0.32 & 0.27 & 0.16 & 0.22 \\
DR & 0.48 & 0.22 & 0.31 & 0.07 & 0.40 & 0.32 & 0.02 & 0.16 & 0.22 & 0.01 & 0.00 & 0.11 & 0.21 & 0.25 & 0.20 \\
BT & 0.34 & 0.14 & 0.82 & 0.00 & 0.93 & 0.58 & 0.88 & 0.63 & 0.83 & 0.07 & 0.08 & 0.20 & 0.43 & 0.38 & 0.45 \\
BI & 0.00 & 0.01 & 0.00 & 0.00 & 0.00 & 0.00 & 0.12 & 0.00 & 0.41 & 0.01 & 0.03 & 0.00 & 0.00 & 0.00 & 0.04 \\
BR & 0.01 & 0.02 & 0.00 & 0.00 & 0.22 & 0.33 & 0.13 & 0.00 & 0.06 & 0.03 & 0.00 & 0.04 & 0.01 & 0.00 & 0.06 \\
TI & 0.00 & 0.01 & 0.07 & 0.00 & 0.38 & 0.05 & 0.34 & 0.59 & 0.65 & 0.04 & 0.00 & 0.01 & 0.01 & 0.01 & 0.15 \\
TR & 0.22 & 0.06 & 0.53 & 0.00 & 0.93 & 0.51 & 0.76 & 0.45 & 0.24 & 0.05 & 0.00 & 0.10 & 0.12 & 0.33 & 0.31 \\
IR & 0.00 & 0.01 & 0.00 & 0.00 & 0.00 & 0.00 & 0.03 & 0.07 & 0.20 & 0.03 & 0.00 & 0.10 & 0.01 & 0.04 & 0.03 \\
\midrule
DBT & 0.61 & 0.27 & 0.80 & 0.00 & 0.94 & 0.27 & 0.81 & 0.55 & 0.74 & 0.10 & 0.03 & 0.49 & 0.50 & 0.77 & 0.49 \\
DBI & 0.23 & 0.23 & 0.53 & 0.34 & 0.30 & 0.04 & 0.16 & 0.31 & 0.44 & 0.02 & 0.01 & 0.50 & 0.25 & 0.10 & 0.25 \\
DBR & 0.30 & 0.19 & 0.41 & 0.01 & 0.61 & 0.27 & 0.16 & 0.02 & 0.35 & 0.02 & 0.00 & 0.32 & 0.13 & 0.37 & 0.22 \\
DTI & 0.02 & 0.08 & 0.38 & 0.00 & 0.34 & 0.04 & 0.41 & 0.48 & 0.61 & 0.07 & 0.03 & 0.30 & 0.12 & 0.04 & 0.21 \\
DTR & 0.21 & 0.11 & 0.36 & 0.00 & 0.69 & 0.27 & 0.76 & 0.28 & 0.44 & 0.09 & 0.00 & 0.14 & 0.12 & 0.34 & 0.27 \\
DIR & 0.22 & 0.10 & 0.18 & 0.03 & 0.20 & 0.03 & 0.05 & 0.11 & 0.23 & 0.00 & 0.00 & 0.15 & 0.01 & 0.02 & 0.09 \\
BTI & 0.00 & 0.02 & 0.06 & 0.00 & 0.13 & 0.01 & 0.32 & 0.58 & 0.72 & 0.03 & 0.00 & 0.00 & 0.00 & 0.00 & 0.13 \\
BTR & 0.01 & 0.03 & 0.47 & 0.00 & 0.81 & 0.41 & 0.64 & 0.30 & 0.32 & 0.04 & 0.00 & 0.01 & 0.01 & 0.03 & 0.22 \\
BIR & 0.00 & 0.01 & 0.00 & 0.00 & 0.00 & 0.00 & 0.00 & 0.00 & 0.03 & 0.01 & 0.00 & 0.00 & 0.00 & 0.00 & 0.00 \\
TIR & 0.00 & 0.01 & 0.02 & 0.00 & 0.09 & 0.03 & 0.13 & 0.34 & 0.21 & 0.05 & 0.00 & 0.00 & 0.00 & 0.00 & 0.06 \\
\midrule
DBTI & 0.02 & 0.09 & 0.48 & 0.00 & 0.15 & 0.00 & 0.29 & 0.44 & 0.58 & 0.07 & 0.04 & 0.25 & 0.04 & 0.01 & 0.18 \\
DBTR & 0.15 & 0.09 & 0.47 & 0.00 & 0.73 & 0.23 & 0.63 & 0.20 & 0.43 & 0.10 & 0.00 & 0.11 & 0.03 & 0.30 & 0.25 \\
DBIR & 0.09 & 0.09 & 0.33 & 0.00 & 0.11 & 0.01 & 0.03 & 0.02 & 0.29 & 0.00 & 0.00 & 0.10 & 0.00 & 0.00 & 0.08 \\
DTIR & 0.01 & 0.02 & 0.24 & -- & 0.28 & 0.04 & 0.30 & 0.24 & 0.39 & 0.06 & 0.00 & 0.04 & 0.01 & 0.01 & 0.13 \\
BTIR & 0.00 & 0.01 & 0.02 & 0.00 & 0.01 & 0.00 & 0.02 & 0.32 & 0.27 & 0.01 & 0.00 & 0.00 & 0.00 & 0.00 & 0.05 \\
\midrule
DBTIR & 0.00 & 0.03 & 0.26 & -- & 0.07 & 0.00 & 0.17 & 0.20 & 0.36 & 0.06 & 0.00 & 0.02 & 0.00 & 0.00 & 0.09 \\
\midrule
E7 raw+F & 0.46 & 0.30 & 0.88 & 0.17 & 0.87 & 0.32 & 0.03 & 0.42 & 0.37 & 0.02 & 0.16 & 0.13 & 0.07 & 0.17 & 0.31 \\
CC-20 & 0.00 & 0.09 & 0.09 & 0.00 & 0.77 & 0.14 & 0.53 & 0.40 & 0.27 & 0.00 & 0.00 & 0.00 & 0.00 & 0.00 & 0.16 \\
CC & 0.42 & 0.34 & 0.69 & 0.04 & 0.75 & 0.29 & 0.92 & 0.73 & 0.24 & 0.17 & 0.13 & 0.26 & 0.18 & 0.38 & 0.40 \\
FV-20 & 0.02 & 0.18 & 0.44 & 0.64 & 0.88 & 0.00 & 0.92 & 0.82 & 0.64 & 0.01 & 0.01 & 0.00 & 0.00 & 0.01 & 0.33 \\
FV-$N_t$ & 0.49 & 0.55 & 0.85 & 0.85 & 0.94 & 0.47 & 0.93 & 0.88 & 0.91 & 0.70 & 0.19 & 0.21 & 0.02 & 0.25 & 0.59 \\
\bottomrule
\end{tabular}
}
\caption{Check 1, task accuracy at the cross-fitted $\alpha$, every set and task, mean of the two folds. Blank: empty set.}
\label{tab:all1}
\end{table}

\begin{table}[H]\centering\scriptsize
\resizebox{\linewidth}{!}{% Generated by scripts/e8_tables.py; do not edit.
\begin{tabular}{lrrrrrrrrrrrrrrr}
\toprule
set & ant. & syn. & sg--pl & pr--pa & capital & currency & landm. & park & product & next & prev & En--Fr & En--De & En--Es & mean \\
\midrule
D & 0.29 & 0.11 & 0.73 & 0.68 & 0.52 & 0.37 & 0.04 & 0.14 & 0.11 & 0.02 & 0.03 & 0.17 & 0.27 & 0.21 & 0.26 \\
B & 0.04 & 0.02 & 0.21 & 0.12 & 0.27 & 0.19 & 0.19 & 0.45 & 0.18 & 0.11 & 0.09 & 0.15 & 0.15 & 0.24 & 0.17 \\
T & 0.37 & 0.07 & 0.76 & 0.75 & 0.81 & 0.29 & 0.31 & 0.38 & 0.30 & 0.44 & 0.17 & 0.35 & 0.30 & 0.35 & 0.40 \\
I & 0.00 & 0.01 & 0.00 & 0.01 & 0.00 & 0.10 & 0.03 & 0.48 & 0.14 & 0.02 & 0.05 & 0.03 & 0.01 & 0.03 & 0.06 \\
R & 0.00 & 0.00 & 0.00 & 0.00 & 0.01 & 0.25 & 0.09 & 0.04 & 0.13 & 0.02 & 0.12 & 0.05 & 0.06 & 0.10 & 0.06 \\
\midrule
DB & 0.33 & 0.16 & 0.79 & 0.69 & 0.68 & 0.22 & 0.22 & 0.22 & 0.16 & 0.12 & 0.05 & 0.25 & 0.27 & 0.29 & 0.32 \\
DT & 0.39 & 0.16 & 0.76 & 0.46 & 0.81 & 0.17 & 0.33 & 0.26 & 0.23 & 0.35 & 0.10 & 0.32 & 0.27 & 0.33 & 0.35 \\
DI & 0.20 & 0.09 & 0.71 & 0.67 & 0.54 & 0.26 & 0.03 & 0.20 & 0.12 & 0.04 & 0.06 & 0.23 & 0.25 & 0.22 & 0.26 \\
DR & 0.30 & 0.11 & 0.73 & 0.36 & 0.65 & 0.36 & 0.03 & 0.09 & 0.13 & 0.03 & 0.04 & 0.16 & 0.26 & 0.24 & 0.25 \\
BT & 0.35 & 0.10 & 0.76 & 0.74 & 0.76 & 0.25 & 0.33 & 0.32 & 0.28 & 0.50 & 0.19 & 0.22 & 0.24 & 0.28 & 0.38 \\
BI & 0.00 & 0.00 & 0.02 & 0.07 & 0.21 & 0.22 & 0.12 & 0.41 & 0.19 & 0.07 & 0.05 & 0.12 & 0.13 & 0.16 & 0.13 \\
BR & 0.03 & 0.01 & 0.22 & 0.01 & 0.55 & 0.17 & 0.18 & 0.37 & 0.19 & 0.10 & 0.05 & 0.14 & 0.15 & 0.23 & 0.17 \\
TI & 0.07 & 0.03 & 0.53 & 0.68 & 0.43 & 0.24 & 0.25 & 0.37 & 0.30 & 0.35 & 0.14 & 0.14 & 0.13 & 0.13 & 0.27 \\
TR & 0.31 & 0.06 & 0.76 & 0.26 & 0.81 & 0.34 & 0.29 & 0.28 & 0.30 & 0.42 & 0.07 & 0.36 & 0.33 & 0.34 & 0.35 \\
IR & 0.01 & 0.02 & 0.00 & 0.02 & 0.03 & 0.06 & 0.01 & 0.44 & 0.12 & 0.02 & 0.07 & 0.05 & 0.07 & 0.10 & 0.07 \\
\midrule
DBT & 0.36 & 0.14 & 0.75 & 0.46 & 0.71 & 0.16 & 0.34 & 0.26 & 0.22 & 0.55 & 0.10 & 0.29 & 0.24 & 0.30 & 0.35 \\
DBI & 0.24 & 0.14 & 0.58 & 0.51 & 0.35 & 0.20 & 0.17 & 0.18 & 0.16 & 0.08 & 0.06 & 0.25 & 0.25 & 0.18 & 0.24 \\
DBR & 0.32 & 0.12 & 0.74 & 0.43 & 0.55 & 0.23 & 0.18 & 0.19 & 0.17 & 0.09 & 0.04 & 0.25 & 0.28 & 0.28 & 0.28 \\
DTI & 0.21 & 0.13 & 0.53 & 0.49 & 0.67 & 0.16 & 0.27 & 0.23 & 0.25 & 0.37 & 0.07 & 0.21 & 0.22 & 0.21 & 0.29 \\
DTR & 0.36 & 0.14 & 0.75 & 0.11 & 0.70 & 0.19 & 0.33 & 0.19 & 0.23 & 0.30 & 0.09 & 0.28 & 0.24 & 0.30 & 0.30 \\
DIR & 0.19 & 0.07 & 0.68 & 0.42 & 0.52 & 0.24 & 0.02 & 0.22 & 0.12 & 0.03 & 0.06 & 0.21 & 0.19 & 0.20 & 0.23 \\
BTI & 0.06 & 0.04 & 0.59 & 0.67 & 0.36 & 0.17 & 0.25 & 0.29 & 0.30 & 0.45 & 0.13 & 0.12 & 0.12 & 0.12 & 0.26 \\
BTR & 0.23 & 0.08 & 0.69 & 0.26 & 0.77 & 0.28 & 0.34 & 0.23 & 0.29 & 0.48 & 0.12 & 0.20 & 0.26 & 0.25 & 0.32 \\
BIR & 0.00 & 0.01 & 0.00 & 0.01 & 0.36 & 0.16 & 0.13 & 0.35 & 0.17 & 0.08 & 0.06 & 0.12 & 0.10 & 0.16 & 0.12 \\
TIR & 0.11 & 0.04 & 0.59 & 0.16 & 0.48 & 0.22 & 0.19 & 0.31 & 0.28 & 0.33 & 0.07 & 0.17 & 0.19 & 0.17 & 0.24 \\
\midrule
DBTI & 0.23 & 0.13 & 0.56 & 0.49 & 0.66 & 0.14 & 0.26 & 0.25 & 0.23 & 0.51 & 0.10 & 0.21 & 0.24 & 0.20 & 0.30 \\
DBTR & 0.34 & 0.14 & 0.72 & 0.11 & 0.71 & 0.13 & 0.33 & 0.19 & 0.22 & 0.53 & 0.10 & 0.23 & 0.26 & 0.27 & 0.31 \\
DBIR & 0.22 & 0.10 & 0.56 & 0.35 & 0.39 & 0.20 & 0.13 & 0.18 & 0.17 & 0.08 & 0.05 & 0.22 & 0.22 & 0.17 & 0.22 \\
DTIR & 0.23 & 0.10 & 0.50 & -- & 0.69 & 0.16 & 0.24 & 0.18 & 0.23 & 0.35 & 0.11 & 0.23 & 0.24 & 0.21 & 0.27 \\
BTIR & 0.08 & 0.05 & 0.60 & 0.16 & 0.32 & 0.14 & 0.18 & 0.23 & 0.29 & 0.47 & 0.12 & 0.16 & 0.12 & 0.15 & 0.22 \\
\midrule
DBTIR & 0.23 & 0.10 & 0.52 & -- & 0.69 & 0.12 & 0.20 & 0.20 & 0.22 & 0.46 & 0.12 & 0.24 & 0.22 & 0.18 & 0.27 \\
\midrule
E7 raw+F & 0.34 & 0.16 & 0.77 & 0.78 & 0.74 & 0.12 & 0.19 & 0.38 & 0.21 & 0.25 & 0.12 & 0.22 & 0.27 & 0.25 & 0.34 \\
CC-20 & 0.31 & 0.13 & 0.68 & 0.57 & 0.73 & 0.12 & 0.32 & 0.21 & 0.13 & 0.08 & 0.09 & 0.12 & 0.13 & 0.14 & 0.27 \\
CC & 0.35 & 0.17 & 0.77 & 0.64 & 0.77 & 0.23 & 0.32 & 0.33 & 0.16 & 0.34 & 0.09 & 0.16 & 0.21 & 0.24 & 0.34 \\
FV-20 & 0.30 & 0.11 & 0.64 & 0.89 & 0.67 & 0.19 & 0.33 & 0.33 & 0.20 & 0.23 & 0.17 & 0.26 & 0.21 & 0.20 & 0.34 \\
FV-$N_t$ & 0.35 & 0.18 & 0.65 & 0.82 & 0.74 & 0.17 & 0.32 & 0.38 & 0.22 & 0.61 & 0.32 & 0.36 & 0.17 & 0.22 & 0.39 \\
\bottomrule
\end{tabular}
}
\caption{Check 2, as Table~\ref{tab:all1}.}
\label{tab:all2}
\end{table}

\end{document}
